\scanpage{88}
\deleted{\unsure{(iii)}} \deleted{\unsure{(iv)|(v)}} i.e. la forme induite sur $M$ \inserted{par une forme ample} est non dégénérée, donc $N=M\oplus M^\perp$, et la forme sur $N$ est la somme des formes induites sur $M,M^\perp$.

(v) La somme directe de deux formes \deleted{positives \unsure{ou}} pseudoamples (amples) est itou.

(vi) Le produit tensoriel de deux formes amples est ample.

(vii) Si $p=0$, de sorte que \unsure{essentiellt|essentielt|essentiellement} $M$ est défini par un vectoriel de dim finie sur $\mathbb{Q}$, alors les formes amples sont les formes \uline{définies positives}. \deleted{\unsure{Plus généralement|Plus généralemt} pour $M. (p)$, $M$ un vectoriel de dim finie sur $\mathbb{Q}$ ; utilisant l'isom canon}

(viii) Cas de $\mathbb{Q}(p)$ : on a $\unsure{\Phi|\mathscr{L}|\mathcal{L}}^s(\mathbb{Q}(p))\simeq\mathbb{Q}$, et on veut que les formes amples correspondent aux nbs \inserted{\uline{\unsure{strictt|strict\textsuperscript{t}|strictement}}} \uline{positifs}.
% archive sheet number: 87
